\section*{Capítulo \ref{ch:b1:lipids} --- Lipídios}

\begin{solution}{exo:b1:lipids:1}
Um \omterm{def:b1:lipids:fattyacid}{lipídio} é uma molécula biológica insolúvel em água e solúvel em
solventes apolares. \omterm{def:b1:lipids:triglyceride}{Gorduras}, \omterm{def:b1:lipids:amphiphilic}{fosfolipídios}, \omterm{def:b1:lipids:amphiphilic}{esteróis} e carotenoides não
compartilham nenhum esqueleto comum, apenas esse comportamento diante da
água, que é o que lhes dá os papéis comuns (estoques, películas,
membranas).
\end{solution}

\begin{solution}{exo:b1:lipids:2}
C18:2 $\Delta^{9,12}$. A última ligação dupla começa no carbono 12 a
partir da carboxila, ou seja, no carbono 6 a partir da extremidade
metila: ômega-6.
\end{solution}

\begin{solution}{exo:b1:lipids:3}
Um \omterm{def:b1:lipids:triglyceride}{triglicerídeo} não tem cabeça \omterm{def:b1:water-small-molecules:water}{polar}: inteiramente hidrofóbico, ele é
excluído da água como uma fase à parte, uma gotícula. Um \omterm{def:b1:lipids:amphiphilic}{fosfolipídio} é
anfifílico: a cabeça \omterm{def:b1:water-small-molecules:water}{polar} dele precisa ficar na água enquanto as caudas
precisam sair dela, o que só uma lâmina de duas moléculas de espessura
consegue satisfazer.
\end{solution}

\begin{solution}{exo:b1:lipids:4}
Cerca de \qty{41}{\celsius}, \qty{-5}{\celsius} e \qty{20}{\celsius} (a
curva do \omterm{def:b1:lipids:amphiphilic}{colesterol} não tem transição brusca; metade da subida dela fica
perto de \qty{20}{\celsius}). Uma bactéria a \qty{10}{\celsius} precisa
da composição insaturada, fluida bem abaixo da temperatura de cultivo.
\end{solution}

\begin{solution}{exo:b1:lipids:5}
$15\,000\times 38 = \qty{570}{MJ}$; $570/8 = 71$ dias. Glicogênio: seco
$570\,000/17 = \qty{33.5}{kg}$, hidratado \qty{134}{kg}.
\end{solution}

\begin{solution}{exo:b1:lipids:6}
C22:0 $>$ C18:0 $>$ C14:0 $>$ C18:1 $>$ C18:3. Entre os ácidos
saturados, as cadeias mais longas fazem mais contatos \omterm{def:b1:water-small-molecules:bonds}{de van der Waals};
uma ligação dupla \emph{cis} elimina mais empacotamento do que quatro
carbonos a mais acrescentam (o C18:1 funde abaixo do C14:0); cada
ligação adicional o baixa outra vez.
\end{solution}

\begin{solution}{exo:b1:lipids:7}
$2\times\qty{140}{\micro m^2}/\qty{0.6}{nm^2} = 2\times 1.4\times
10^{-10}/6\times 10^{-19} = \num{4.7e8}$ moléculas; massa
$\num{4.7e8}\times 750/\num{6e23} = \qty{5.8e-13}{g} = \qty{0.58}{pg}$.
\end{solution}

\begin{solution}{exo:b1:lipids:8}
Volume \qty{1.11}{mL} $= \qty{1.11e-6}{m^3}$; superfície das esferas
$= 6V/d = 6\times 1.11\times 10^{-6}/10^{-6} = \qty{6.7}{m^2}$. Uma
única gota desse volume: $d = \qty{1.28}{cm}$, superfície
\qty{5.2}{cm^2}: a emulsão tem treze mil vezes mais interface. A lipase
é solúvel em água e só age na interface; a velocidade é proporcional a
ela.
\end{solution}

\begin{solution}{exo:b1:lipids:9}
Acima de $T_m$ os anéis rígidos dele atrapalham o movimento das cadeias
vizinhas e reduzem a \omterm{def:b1:lipids:fluidity}{fluidez}; abaixo de $T_m$ o volume dele impede que
as cadeias se empacotem no gel ordenado e impede o congelamento. Sem
empacotamento cooperativo possível, não há temperatura em que toda a
bicamada mude de estado de uma vez: a transição se espalha.
\end{solution}

\begin{solution}{exo:b1:lipids:10}
Cadeias todas saturadas: a \qty{25}{\celsius} as membranas já estão
perto da transição delas e rígidas; a \qty{5}{\celsius} elas
gelificam. Membranas gelificadas deixam as proteínas de parar de
funcionar e, ao reaquecer ou sob esforço mecânico, rachar e vazar: as
\omterm{def:b1:cell-unit-of-life:cell}{células} perdem o conteúdo e o \omterm{def:b1:body-plans-tissues:tissue}{tecido} morre (lesão pelo frio). O tipo
selvagem dessatura os \omterm{def:b1:lipids:fattyacid}{lipídios} dele no outono e continua fluido a
\qty{5}{\celsius}.
\end{solution}

\begin{solution}{exo:b1:lipids:11}
\qty{1.07}{kg} de água por quilograma de \omterm{def:b1:lipids:triglyceride}{gordura}. Oxigênio:
\qty{2000}{L}; com \qty{5}{\%} de extração de um ar com \qty{21}{\%} de
oxigênio, o ar necessário é $2000/(0.21\times 0.05) = \qty{190}{m^3}$;
água perdida a \qty{30}{g/m^3}: \qty{5.7}{kg}. O camelo perde cinco
vezes mais água respirando do que a \omterm{def:b1:lipids:triglyceride}{gordura} rende: a corcova é um
estoque de energia, não de água.
\end{solution}

\begin{solution}{exo:b1:lipids:12}
Dada uma bicamada, o \omterm{prop:b1:water-small-molecules:properties}{efeito hidrofóbico} faz com que ela se refeche, se
feche em vesículas e aceite \omterm{def:b1:lipids:fattyacid}{lipídios} novos sem nenhum aporte de energia:
a montagem é espontânea. Mas a \omterm{def:b1:cell-unit-of-life:cell}{célula} fabrica os \omterm{def:b1:lipids:fattyacid}{lipídios} dela com
enzimas alojadas numa membrana já existente (o RE), que os inserem no
lugar; uma membrana nova cresce pela expansão de uma antiga e é
repartida na divisão, de modo que toda membrana de toda \omterm{def:b1:cell-unit-of-life:cell}{célula} descende
das membranas da \omterm{def:b1:cell-unit-of-life:cell}{célula} anterior. A automontagem explica a física da
lâmina, não a origem dela na \omterm{def:b1:cell-unit-of-life:cell}{célula}.
\end{solution}

\begin{solution}{pb:b1:lipids:1}
\textbf{1.} C16:0 (0), C18:1 $\Delta^9$ (1), C22:6 $\Delta^{4,7,10,13,16,19}$
(6).
\textbf{2.} Água morna: $0.40\times 0 + 0.45\times 1 + 0.15\times 6 = 1.35$;
água fria: $0 + 0.40 + 2.10 = 2.50$ ligações duplas por cadeia.
\textbf{3.} Cada ligação \emph{cis} dobra a cadeia e impede o
empacotamento próximo com as vizinhas; menos contatos \omterm{def:b1:water-small-molecules:bonds}{de van der Waals},
menos calor necessário para desordenar as cadeias, $T_m$ menor.
\textbf{4.} $1.15$ ligação a mais por cadeia: cerca de
\qty{45}{\celsius} a menos.
\textbf{5.} Num gel os \omterm{def:b1:lipids:fattyacid}{lipídios} não conseguem se mover: os carreadores
não conseguem mudar de conformação, as bombas param, os \omterm{def:b1:membranes-transport:transporters}{canais} travam;
os gradientes iônicos se esgotam, a condução nervosa falha, e o peixe
fica paralisado e morre do frio que um peixe aclimatado tolera.
\textbf{6.} As dessaturases, que introduzem ligações duplas nas cadeias
de acila graxa; elas são proteínas \omterm{def:b1:membranes-transport:proteins}{integrais} do \omterm{def:b1:eukaryotic-cell:endomembrane}{retículo
endoplasmático}, onde os \omterm{def:b1:lipids:fattyacid}{lipídios} são fabricados.
\textbf{7.} Nos dois casos ele tampona a \omterm{def:b1:lipids:fluidity}{fluidez}: enrijece a membrana
fria, muito insaturada, acima do baixo $T_m$ dela, e impede que a
membrana morna gelifique num dia fresco.
\textbf{8.} $30\,000\times 38 = \qty{1140}{MJ}$.
\textbf{9.} $1140/60 = 19$ dias.
\textbf{10.} Seco: $\num{1.14e6}/17 = \qty{67}{kg}$; hidratado:
\qty{268}{kg}.
\textbf{11.} $268 - 30 = \qty{238}{kg}$, quase metade da massa do
camelo.
\textbf{12.} Uma camada de \omterm{def:b1:lipids:triglyceride}{gordura} sob toda a pele isolaria o corpo e o
impediria de dissipar calor no deserto; concentrar a \omterm{def:b1:lipids:triglyceride}{gordura} em uma
única corcova deixa o resto da pele livre para perder calor.
\textbf{13.} $1.07\times 30 = \qty{32}{kg}$ de água.
\textbf{14.} $2.0\times 30\,000 = \qty{60000}{L} = \qty{60}{m^3}$ de
oxigênio.
\textbf{15.} $60/(0.21\times 0.05) = \qty{5700}{m^3}$ de ar.
\textbf{16.} $5700\times 30 = \qty{171}{kg}$ de água.
\textbf{17.} Perde cinco vezes o que ganha: oxidar a corcova custa
água. A corcova é um estoque de energia que deixa o camelo passar sem
comida; a economia de água dele vem dos rins, da tolerância à
desidratação e da capacidade de deixar a temperatura subir.
\textbf{18.} $500\times 3.5\times 6 = \qty{10.5}{MJ}$ armazenados como calor e liberados à noite; evaporá-los teria custado $10\,500/2.4 = \qty{4.4}{kg}$ de água por dia.
\textbf{19.} Uma membrana só consegue esticar alguns por cento, de modo
que dobrar o volume de uma hemácia exige um grande excesso de membrana
dobrado na forma bicôncava em repouso, e um \omterm{def:b1:eukaryotic-cell:cytoskeleton}{citoesqueleto} que a deixe
desdobrar sem rasgar: as hemácias do camelo são ovais, com mais membrana
por volume que as humanas.
\textbf{20.} $2\times 80\times 10^{-12}/0.65\times 10^{-18} = \num{2.5e8}$.
\textbf{21.} $\num{2.5e8}\times\num{8e12}\times 40 = \num{7.9e22}$.
\textbf{22.} $\num{7.9e22}\times 750/\num{6e23} = \qty{99}{g}$.
\textbf{23.} A difusão lateral mantém a cabeça na água e as caudas no
\omterm{def:b1:lipids:triglyceride}{óleo} a cada passo, sem custo nenhum; passar de um folheto ao outro exige
que a cabeça \omterm{def:b1:water-small-molecules:water}{polar} atravesse o núcleo hidrofóbico, uma barreira
energética grande que a agitação térmica vence uma vez por semana.
\textbf{24.} Uma diferença entre folhetos só se desfaz na velocidade com
que os \omterm{def:b1:lipids:fattyacid}{lipídios} passam de um ao outro; a uma vez por semana e por
molécula, uma assimetria estabelecida por enzimas que transferem
\omterm{def:b1:lipids:fattyacid}{lipídios} específicos (as flipases) persiste indefinidamente contra a
passagem espontânea.
\textbf{25.} Cerca de \qty{240}{kg} poupados --- quase metade da massa
corporal dele --- por carregar \qty{30}{kg} de \omterm{def:b1:lipids:triglyceride}{gordura} em vez de
\qty{270}{kg} de glicogênio hidratado; e a corcova não é um estoque de
água: oxidá-la perde na respiração cinco vezes mais água do que produz.
\end{solution}
